%%=============================================================================
%% Samenvatting
%%=============================================================================

% TODO: De "abstract" of samenvatting is een kernachtige (~ 1 blz. voor een
% thesis) synthese van het document.
%
% Een goede abstract biedt een kernachtig antwoord op volgende vragen:
%
% 1. Waarover gaat de bachelorproef?
% 2. Waarom heb je er over geschreven?
% 3. Hoe heb je het onderzoek uitgevoerd?
% 4. Wat waren de resultaten? Wat blijkt uit je onderzoek?
% 5. Wat betekenen je resultaten? Wat is de relevantie voor het werkveld?
%
% Daarom bestaat een abstract uit volgende componenten:
%
% - inleiding + kaderen thema
% - probleemstelling
% - (centrale) onderzoeksvraag
% - onderzoeksdoelstelling
% - methodologie
% - resultaten (beperk tot de belangrijkste, relevant voor de onderzoeksvraag)
% - conclusies, aanbevelingen, beperkingen
%
% LET OP! Een samenvatting is GEEN voorwoord!

%%---------- Nederlandse samenvatting -----------------------------------------
%
% TODO: Als je je bachelorproef in het Engels schrijft, moet je eerst een
% Nederlandse samenvatting invoegen. Haal daarvoor onderstaande code uit
% commentaar.
% Wie zijn bachelorproef in het Nederlands schrijft, kan dit negeren, de inhoud
% wordt niet in het document ingevoegd.

% \IfLanguageName{english}{%
% \selectlanguage{dutch}
% \chapter*{Samenvatting}
% \lipsum[1-4]
% \selectlanguage{english}
% }{}

%%---------- Samenvatting -----------------------------------------------------
% De samenvatting in de hoofdtaal van het document

\chapter*{\IfLanguageName{dutch}{Samenvatting}{Abstract}}

Digitale toegankelijkheid is binnen het hoger onderwijs een noodzakelijke voorwaarde om gelijke participatiekansen te garanderen. De HOGENT-campusapp wordt intensief gebruikt voor communicatie, planning en toegang tot studie-informatie, maar wanneer toegankelijkheidsrichtlijnen onvoldoende zijn geïntegreerd, ontstaan er concrete drempels voor studenten met een visuele, auditieve, motorische of cognitieve beperking. Vanuit die probleemstelling vertrekt deze bachelorproef met de centrale onderzoeksvraag: Hoe kan de HOGENT-campusapp inclusiever worden ontworpen voor studenten met een beperking, rekening houdend met hun specifieke noden op het vlak van digitale toegankelijkheid? Om deze vraag te beantwoorden wordt een gefaseerde en iteratieve methodologie toegepast. Eerst wordt een literatuurstudie uitgevoerd rond WCAG 2.1/2.2, relevante Europese en Belgische regelgeving en gangbare evaluatiemethoden voor mobiele toegankelijkheid. Vervolgens wordt de huidige applicatie geaudit via een combinatie van geautomatiseerde tools (Accessibility Insights, Lighthouse en axe DevTools) en manuele testen, waaronder toetsenbordnavigatie en schermlezercompatibiliteit. De vastgestelde problemen worden geprioriteerd op basis van gebruikersimpact en technische haalbaarheid, waarna gerichte verbeteringen op UI- en componentniveau worden geïmplementeerd en opnieuw geëvalueerd ten opzichte van een nulmeting. Het onderzoek beoogt een aantoonbare verbetering van de toegankelijkheid van de campusapp, met een reductie van kritieke WCAG-gerelateerde problemen zonder significante negatieve impact op performantie. De bijdrage voor het werkveld ligt in een praktisch toepasbaar verbeterkader en concrete aanbevelingen om toegankelijkheid structureel te verankeren in ontwerp-, ontwikkel- en testprocessen van mobiele applicaties. Tegelijk kent het onderzoek beperkingen: geautomatiseerde tools detecteren slechts een deel van alle toegankelijkheidsproblemen, manuele evaluatie is tijdsintensief en gebruikerstesten vereisen voldoende representatieve participatie. Verdere validatie in bredere praktijksituaties blijft daarom aangewezen.


